\documentclass{article}
\usepackage[UTF8]{ctex}
\usepackage{amsmath,amssymb,amsfonts}
\usepackage{enumitem}
\usepackage{geometry}
\geometry{a4paper, margin=1in}
% === 错题本排版优化（不参与题目标识） ===
\setlength{\parindent}{0pt}
\setlength{\parskip}{0.35em plus 0.10em minus 0.08em}
\setlength{\emergencystretch}{3em}
\setlist{itemsep=0.25em, topsep=0.40em, parsep=0pt}
\allowdisplaybreaks[2]
% === 错题本排版优化结束 ===
\begin{document}
\section*{华南师范大学2026年数学分析真题}
\section*{161．华南师范大学2026年数学分析试卷}

\subsection{题目1}
1． \(\lim _{n \rightarrow \infty} \sqrt{n}\left(\sqrt{n^{2}+1}-\sqrt{n+1}\right)\) ．

\subsection{题目2}
2． \(\lim _{n \rightarrow \infty} n(\sqrt[n]{a}-1)(a>0)\) ．

\subsection{题目3}
3． \(\lim _{x \rightarrow 0} \frac{\sqrt{\cos x}-\sqrt{1+x \arctan x}}{2 x^{2}}\) ．

\subsection{题目4}
4． \(\lim _{x \rightarrow+\infty} e^{-x}\left(1+\frac{1}{x}\right)^{x^{2}}\)

\subsection{题目5}
5． \(\lim _{(x, y) \rightarrow(+\infty,+\infty)}\left(1+\frac{1}{x y}\right)^{x \sin y}\)

\subsection{题目6}
1．计算定积分 \(\int_{0}^{100 \pi} \sqrt{1-\cos 2 x} \mathrm{~d} x\) ．

\subsection{题目7}
2．计算 \(\iint_{D} \sqrt{1-\sin ^{2}(x+y)} \mathrm{d} x \mathrm{~d} y\) ，其中 \(D=\left\{(x, y) \mid 0 \leq y \leq x, 0 \leq x \leq \frac{\pi}{2}\right\}\) ．

\subsection{题目8}
3．已知 \(\left\{\begin{array}{l}x=t^{2}+t, \\ y=\ln (1+t)\end{array}\right.\) ．求 \(\frac{\mathrm{d} y}{\mathrm{~d} x}\) 与 \(\frac{\mathrm{d}^{2} y}{\mathrm{~d} x^{2}}\) ．

\subsection{题目9}
4．已知 \(u=f\left(x y, \frac{y}{x}\right)\) ，其中 \(f\) 二阶连续可微，求 \(\frac{\partial u}{\partial x}\) 与 \(\frac{\partial^{2} u}{\partial x \partial y}\) ．

三、讨论题．10分．\\
设二元函数

\[
f(x, y)= \begin{cases}\left(x^{2}+y^{2}\right)^{\alpha} \sin \frac{1}{x^{2}+y^{2}}, & x^{2}+y^{2} \neq 0 \\ 0, & x^{2}+y^{2}=0\end{cases}
\]

分别讨论 \(f(x, y)\) 在 \((0,0)\) 处连续，偏导数存在时，\(\alpha\) 的取值范围．

\subsection{题目10}
1．求 \(\sum_{n=1}^{\infty} \frac{1}{n \cdot 2^{n}}\) 的和．

\subsection{题目11}
2．设 \(f_{n}(x)=n^{\alpha} x e^{-n x}\) ，求 \(\alpha\) 的取值范围，使得函数列 \(\left\{f_{n}(x)\right\}\) 在［0，1］上一致收敛．

\subsection{题目12}
3．判断级数或反常积分的敛散性．

\begin{itemize}
  \item[（1）] \(\sum_{n=1}^{\infty} \frac{(-1)^{n+1}}{n} \frac{a}{1+a^{n}}(a>0)\) ．
  \item[（2）] \(\int_{0}^{\frac{1}{e}} \frac{1}{x^{p} \ln x} \mathrm{~d} x(p>0)\) ．
\end{itemize}

五．证明题．10 分．\\
已知函数 \(f(x)\) 在 \((a, b)\) 内二阶可导，\(f^{\prime}(x)\) 在 \([a, b]\) 上连续，且满足 \(f(a)=f(b)=0, \int_{a}^{b} f(x) \mathrm{d} x=0\) ，证\\
明：存在 \(\xi \in(a, b)\) ，使得 \(f^{\prime \prime}(\xi)=f(\xi)\) ．
\end{document}
